\documentclass[10pt,a4paper]{amsart}
\usepackage{amsmath,amssymb,amsthm,graphicx}
\usepackage{verbatim}
\newenvironment{reference}{\comment}{\endcomment}
%\newenvironment{reference}{}{}
\newenvironment{slogan}{\comment}{\endcomment}
\newenvironment{history}{\comment}{\endcomment}
% For commutative diagrams we use Xy-pic
\usepackage[all]{xy}

% We use 2cell for 2-commutative diagrams.
\xyoption{2cell}
% Theorem environments.
%
\theoremstyle{plain}
\newtheorem{theorem}[subsection]{Theorem}
\newtheorem{proposition}[subsection]{Proposition}
\newtheorem{lemma}[subsection]{Lemma}

\theoremstyle{definition}
\newtheorem{definition}[subsection]{Definition}
\newtheorem{example}[subsection]{Example}
\newtheorem{exercise}[subsection]{Exercise}
\newtheorem{situation}[subsection]{Situation}

\theoremstyle{remark}
\newtheorem{remark}[subsection]{Remark}
\newtheorem{remarks}[subsection]{Remarks}

\numberwithin{equation}{subsection}

% Macros
%
\def\lim{\mathop{\mathrm{lim}}\nolimits}
\def\colim{\mathop{\mathrm{colim}}\nolimits}
\def\Spec{\mathop{\mathrm{Spec}}}
\def\Hom{\mathop{\mathrm{Hom}}\nolimits}
\def\Ext{\mathop{\mathrm{Ext}}\nolimits}
\def\SheafHom{\mathop{\mathcal{H}\!\mathit{om}}\nolimits}
\def\SheafExt{\mathop{\mathcal{E}\!\mathit{xt}}\nolimits}
\def\Sch{\mathit{Sch}}
\def\Mor{\mathop{\mathrm{Mor}}\nolimits}
\def\Ob{\mathop{\mathrm{Ob}}\nolimits}
\def\Sh{\mathop{\mathit{Sh}}\nolimits}
\def\NL{\mathop{N\!L}\nolimits}
\def\CH{\mathop{\mathrm{CH}}\nolimits}
\def\proetale{{pro\text{-}\acute{e}tale}}
\def\etale{{\acute{e}tale}}
\def\QCoh{\mathit{QCoh}}
\def\Ker{\mathop{\mathrm{Ker}}}
\def\Im{\mathop{\mathrm{Im}}}
\def\Coker{\mathop{\mathrm{Coker}}}
\def\Coim{\mathop{\mathrm{Coim}}}

% Boxtimes
%
\DeclareMathSymbol{\boxtimes}{\mathbin}{AMSa}{"02}

%
% Macros for moduli stacks/spaces
%
\def\QCohstack{\mathcal{QC}\!\mathit{oh}}
\def\Cohstack{\mathcal{C}\!\mathit{oh}}
\def\Spacesstack{\mathcal{S}\!\mathit{paces}}
\def\Quotfunctor{\mathrm{Quot}}
\def\Hilbfunctor{\mathrm{Hilb}}
\def\Curvesstack{\mathcal{C}\!\mathit{urves}}
\def\Polarizedstack{\mathcal{P}\!\mathit{olarized}}
\def\Complexesstack{\mathcal{C}\!\mathit{omplexes}}
% \Pic is the operator that assigns to X its picard group, usage \Pic(X)
% \Picardstack_{X/B} denotes the Picard stack of X over B
% \Picardfunctor_{X/B} denotes the Picard functor of X over B
\def\Pic{\mathop{\mathrm{Pic}}\nolimits}
\def\Picardstack{\mathcal{P}\!\mathit{ic}}
\def\Picardfunctor{\mathrm{Pic}}
\def\Deformationcategory{\mathcal{D}\!\mathit{ef}}
\usepackage{cleveref}
\crefname{originallemma}{original lemma}{original lemmas}
\Crefname{originallemma}{Original lemma}{Original lemmas}
\makeatletter
\newlabel{lemma-check-separated-dvr}{{0ARJ}{source}{}{sourcepointer-1}{}}
\newlabel{lemma-check-separated-dvr@cref}{{[originallemma][0][]0ARJ}{[1][0][]source}{}{}{}}
\makeatother
\crefname{originallemma}{original lemma}{original lemmas}
\Crefname{originallemma}{Original lemma}{Original lemmas}
\makeatletter
\newlabel{lemma-weak-chow}{{089J}{source}{}{sourcepointer-2}{}}
\newlabel{lemma-weak-chow@cref}{{[originallemma][0][]089J}{[1][0][]source}{}{}{}}
\makeatother
\crefname{originallemma}{original lemma}{original lemmas}
\Crefname{originallemma}{Original lemma}{Original lemmas}
\makeatletter
\newlabel{limits-lemma-check-universally-closed-Noetherian}{{05JY}{source}{}{sourcepointer-3}{}}
\newlabel{limits-lemma-check-universally-closed-Noetherian@cref}{{[originallemma][0][]05JY}{[1][0][]source}{}{}{}}
\makeatother
\crefname{originallemma}{original lemma}{original lemmas}
\Crefname{originallemma}{Original lemma}{Original lemmas}
\makeatletter
\newlabel{spaces-morphisms-lemma-image-proper-is-proper}{{08AJ}{source}{}{sourcepointer-4}{}}
\newlabel{spaces-morphisms-lemma-image-proper-is-proper@cref}{{[originallemma][0][]08AJ}{[1][0][]source}{}{}{}}
\makeatother
\crefname{originallemma}{original lemma}{original lemmas}
\Crefname{originallemma}{Original lemma}{Original lemmas}
\makeatletter
\newlabel{spaces-morphisms-lemma-universally-closed-local}{{03IT}{source}{}{sourcepointer-5}{}}
\newlabel{spaces-morphisms-lemma-universally-closed-local@cref}{{[originallemma][0][]03IT}{[1][0][]source}{}{}{}}
\makeatother
\title{Existence and uniqueness of lifts along discretely valued diagrams implies properness, by reducing to the weak Chow cover}
\author{The Stacks Project Authors}
\date{stacks-project:spaces-cohomology.tex:local-snapshot-20261012; GNU FDL 1.2 or later}
\begin{document}
\maketitle
\noindent\textit{Editorial note.} All delivered mathematical text is the author's own, copied from the pinned Stacks Project chapter file spaces-cohomology.tex; no statement or proof step was rewritten, reordered or removed.\par\medskip
\noindent\textit{Editorial note.} Preamble provenance: the pinned chapter file opens with \textbackslash{}input\{preamble\} and carries no preamble text of its own, so the project's single Stacks-wide preamble staged at sources/stacks-spaces-cohomology/preamble.tex is excerpted verbatim as three spans, lines 8-12 (verbatim package and the project's own reference/slogan/history display environments), lines 14-18 (Xy-pic setup) and lines 154-217 (theorem declarations and mathematical macros). The wrapper class and the amsmath/amssymb/amsthm/graphicx base packages are supplied by the assembler. No macro was invented or redefined.\par\medskip
\noindent\textit{Editorial note.} Same-chapter pointers are located at their \textbackslash{}label lines inside the pinned chapter file; cross-chapter pointers are located in the project's own tag list staged in the snapshot (sources/stacks-chow/tags.txt), at the verbatim lines carrying their canonical tags.\par\medskip
\noindent\textit{Allowed standard prerequisites (editorial).}\begin{itemize}
\item The separatedness criterion from the same valuative shape, and locality of universal closedness. Same chapter spaces-cohomology.tex, label lemma-check-separated-dvr (invoked at proof line 3710) and Morphisms of Spaces, label spaces-morphisms-lemma-universally-closed-local (invoked at proof line 3713). The first gives separatedness so that properness reduces to universal closedness; the second lets universal closedness be checked etale locally, reducing to a Noetherian affine base.
\item The weak form of Chow's lemma. Same chapter spaces-cohomology.tex, label lemma-weak-chow (invoked at proof line 3716). It produces a proper surjective morphism \$X' \textbackslash{}to X\$ with \$X'\$ a scheme, so that some properties of \$X \textbackslash{}to Y\$ can be tested on \$X'\$.
\item Inheritance of properness along a proper surjective morphism, and the valuative criterion over a Noetherian base. Morphisms of Spaces, label spaces-morphisms-lemma-image-proper-is-proper (invoked at proof line 3718) and Limits, label limits-lemma-check-universally-closed-Noetherian (invoked at proof line 3720). The first shows the claim that \$X' \textbackslash{}to \textbackslash{}Spec(A)\$ is universally closed implies the lemma; the second reduces that claim to lifting dotted arrows for diagrams from the spectrum of a discrete valuation ring.
\item The valuative criterion for proper morphisms of schemes. Standard, used at proof lines 3730-3732. Applied to the proper scheme morphism \$X' \textbackslash{}times\_\{X, b\} \textbackslash{}Spec(A) \textbackslash{}to \textbackslash{}Spec(A)\$, it lifts the dotted arrow \$b\$ to the required arrow \$a\$.
\end{itemize}\medskip
\ifdefined\hypertarget\hypertarget{sourcepointer-1}{}\fi
\noindent\textit{Source reference.} Original lemma 0ARJ: sources/stacks-spaces-cohomology/pinned-original-source.tex, source lines 3632--3632. Navigation pointer for \textbackslash{}ref\{lemma-check-separated-dvr\} used inside the counted span (line 3710); the argument that consumes it is retained in the span.\par\medskip
\ifdefined\hypertarget\hypertarget{sourcepointer-2}{}\fi
\noindent\textit{Source reference.} Original lemma 089J: sources/stacks-spaces-cohomology/pinned-original-source.tex, source lines 3498--3498. Navigation pointer for \textbackslash{}ref\{lemma-weak-chow\} used inside the counted span (line 3716); the argument that consumes it is retained in the span.\par\medskip
\ifdefined\hypertarget\hypertarget{sourcepointer-3}{}\fi
\noindent\textit{Source reference.} Original lemma 05JY: sources/stacks-chow/tags.txt, source lines 7182--7182. Navigation pointer for \textbackslash{}ref\{limits-lemma-check-universally-closed-Noetherian\} used inside the counted span (line 3720); the argument that consumes it is retained in the span.\par\medskip
\ifdefined\hypertarget\hypertarget{sourcepointer-4}{}\fi
\noindent\textit{Source reference.} Original lemma 08AJ: sources/stacks-chow/tags.txt, source lines 10534--10534. Navigation pointer for \textbackslash{}ref\{spaces-morphisms-lemma-image-proper-is-proper\} used inside the counted span (line 3718); the argument that consumes it is retained in the span.\par\medskip
\ifdefined\hypertarget\hypertarget{sourcepointer-5}{}\fi
\noindent\textit{Source reference.} Original lemma 03IT: sources/stacks-chow/tags.txt, source lines 4615--4615. Navigation pointer for \textbackslash{}ref\{spaces-morphisms-lemma-universally-closed-local\} used inside the counted span (line 3713); the argument that consumes it is retained in the span.\par\medskip
\begin{lemma}
\label{lemma-check-proper-dvr}
Let $S$ be a scheme. Let $f : X \to Y$ be a morphism of algebraic
spaces over $S$. Assume
\begin{enumerate}
\item $Y$ is locally Noetherian,
\item $f$ is of finite type and quasi-separated,
\item for every commutative diagram
$$
\xymatrix{
\Spec(K) \ar[r] \ar[d] & X \ar[d] \\
\Spec(A) \ar[r] \ar@{-->}[ru] & Y
}
$$
where $A$ is a discrete valuation ring and $K$ its fraction field,
there is a unique dotted arrow making the diagram commute.
\end{enumerate}
Then $f$ is proper.
\end{lemma}

\begin{proof}
It suffices to prove $f$ is universally closed because $f$ is separated by
Lemma \ref{lemma-check-separated-dvr}.
To do this we may work \'etale locally on $Y$
(Morphisms of Spaces, Lemma
\ref{spaces-morphisms-lemma-universally-closed-local}).
Hence we may assume $Y = \Spec(A)$ is a Noetherian affine scheme.
Choose $X' \to X$ as in the weak form of Chow's lemma
(Lemma \ref{lemma-weak-chow}). We claim that $X' \to \Spec(A)$
is universally closed. The claim implies the lemma by
Morphisms of Spaces, Lemma \ref{spaces-morphisms-lemma-image-proper-is-proper}.
To prove this, according to
Limits, Lemma \ref{limits-lemma-check-universally-closed-Noetherian}
it suffices to prove that in every solid commutative diagram
$$
\xymatrix{
\Spec(K) \ar[r] \ar[d] & X' \ar[r] & X \ar[d] \\
\Spec(A) \ar[rr] \ar@{-->}[ru]^a \ar@{-->}[rru]_b & & Y
}
$$
where $A$ is a dvr with fraction field $K$ we can find the
dotted arrow $a$. By assumption we can find the dotted arrow $b$.
Then the morphism $X' \times_{X, b} \Spec(A) \to \Spec(A)$
is a proper morphism of schemes and by the valuative criterion
for morphisms of schemes we can lift $b$ to the desired morphism $a$.
\end{proof}

\end{document}
